\section{Wire comparisons and global offsets}
\label{sec:offsets}

We now relate the small constraints in $G(\beta)$ to values at complete
binary addresses. A winning game gives sets of possible values, with
one scalar image equality for each comparison. A graph isomorphism
requires more: values must be chosen simultaneously and represented
by sums of functions of adjacent coordinates. This section proves
both statements for arbitrary right-hand sides $\beta$.

\subsection{Wire shifts and image sets}

Fix a comparison $e:P\to R$ and $a\in D_e$. Along its wire, use
the address $a$ at the first endpoint and every internal type, and
$\sigma_e(a)$ at the last endpoint. For a link $\ell$ of this wire,
write $B(i,\ell,a)$ for the square specified by these addresses at
row $i$. Define
\begin{equation}
 \Delta_\ell(a)=\sum_{i=1}^r\beta_{B(i,\ell,a)},
 \qquad
 \Delta_e(a)=\sum_{\ell\text{ on }e}\Delta_\ell(a).
 \label{eq:wire-shift}
\end{equation}
All sums in this section are over $\F$. Each summand belonging to
row $i$ depends only on $(a_i,a_{i+1})$.

\begin{proposition}[Image consistency]
\label{prop:image-consistency}
For the construction of Section~\ref{sec:encoding}, suppose
$G(0)\equiv_r G(\beta)$. There are nonempty sets
$Q_{P,a}\subseteq\F^{d_P}$ for every real type $P$ and every
$a\in\{0,1\}^r$ such that, for every comparison $e:P\to R$,
\begin{equation}
 \lambda_{e,0}(Q_{P,a})
   =\Delta_e(a)+\lambda_{e,1}(Q_{R,\sigma_e(a)})
 \qquad(a\in D_e).
 \label{eq:image-consistency}
\end{equation}
Here the equality concerns scalar image sets for one comparison at
a time, not a simultaneous choice of vectors for all comparisons.
\end{proposition}

The condition has a local-support interpretation, as in classical arc
consistency~\cite[Section~4]{Mackworth1977}: each value at either endpoint
has some supporting value at the other endpoint for the one displayed
comparison. Those supports may differ between comparisons. In
particular, even a self-comparison does not require its two endpoint
values to be the same vector. The proposition and its later converse
are proved here for the stated scalar-image condition.

\begin{proof}
Fix a winning $(r+1)$-pair strategy, supplied by
Lemma~\ref{lem:wl-game}. We will define $Q$ using histories in which
Spoiler selects only zero-valuation vertices of blocks in $G(0)$.
Their responses are valuations of the same blocks in $G(\beta)$:
Lemma~\ref{lem:degree} and the degree intervals preserve every base
color.

We first record a consequence of winning that holds at every such
history. If two marked blocks share a readout, their responses agree
on that readout. Otherwise some shared bit has different values.
Discard all pairs except those two and select the zero vertex of
the shared bit pair. It is adjacent to both source blocks, whereas
no vertex of the corresponding target bit pair is adjacent to both
inconsistent target valuations. This contradicts winning. The test
is a hypothetical continuation: it establishes consistency of the
current responses without being performed in the histories used
below. It needs only three pairs after the discards.

For a type $t$ and a full allowed address $a$ at that type, its
\emph{site ring} consists of the $r$ site blocks with readouts
$x_{i,t,a_i,a_{i+1}}$. When their zero valuations are marked, let
$h_t$ be the sum of the response vectors on these readouts.
Starting with this ring and one spare pair, we can transfer the sum
across any incident link with matching addresses.

For each row, place its square's zero valuation before discarding
the old site pair. This uses at most $r+1$ slots, and consistency
copies the horizontal readout at the departure type. After all rows
are replaced, adjacent square responses agree on their shared
vertical readout. Summing Equation~\eqref{eq:square} around this
square ring cancels every vertical value twice. Thus, for the two
horizontal sums $h_t,h_u$ on a link $\ell=tu$,
\begin{equation}
                  L_{\ell,t}h_t+L_{\ell,u}h_u
                  =\Delta_\ell(a).
 \label{eq:link-sum}
\end{equation}
Now replace each square by its arrival-side site, again placing
before discarding. The arrival sum is copied exactly. The procedure
works in either direction. Along a complete wire, the internal
sums are scalars and cancel between consecutive links. We obtain
the endpoint equality with shift $\Delta_e(a)$.

Define $Q_{P,a}$ to be the set of sums obtained from all
zero-block-only histories ending with exactly the site ring at
$(P,a)$ marked. These sets are nonempty because the ring can be
queried from the empty position. They include histories with every
slot ordering. If desired, an ordering can be changed using the
spare slot: duplicate a marked zero-site vertex before discarding
its former pair. Equality forces the duplicate response to agree,
so successive relocations do not change the ring values.

Continue any history realizing a value in $Q_{P,a}$ along the wire
to obtain one image inclusion in
Equation~\eqref{eq:image-consistency}. Starting from any history
realizing a value in $Q_{R,\sigma_e(a)}$ and traversing backward
gives the reverse inclusion. In the backward traversal the
coordinate injections are inverted only on their images. All
resulting histories still consist of zero-block choices.
No memorylessness or common simultaneous vector choice has been
used. This remains true for a loop: returning to the same ring
may produce a different member of its set.
\end{proof}

\subsection{Lifting sums around a cycle}

The preceding proposition imposes only image conditions. To obtain
an isomorphism, we need actual offsets on every shared readout.
The relevant restriction on full-address values is elementary.

\begin{definition}
For a nonempty product $A=\prod_{i=1}^r A_i$ and a vector space $V$
over $\F$, a function $h:A\to V$ has \emph{row-sum form} if there
are functions $h_i:A_i\times A_{i+1}\to V$ such that
\begin{equation}
                 h(a)=\sum_{i=1}^r h_i(a_i,a_{i+1}).
 \label{eq:row-sum}
\end{equation}
Indices are cyclic, so $A_{r+1}=A_1$ and $a_{r+1}=a_1$.
\end{definition}

Restriction to a product subdomain preserves row-sum form, as do
linear maps on $V$ and pullback by coordinatewise maps of the
domain. In particular, all functions in
Equation~\eqref{eq:wire-shift} have row-sum form.

\begin{lemma}[Cycle lifting]
\label{lem:cycle-lifting}
Let $r\ge3$, let each $A_i$ be nonempty, and let
$g_i:A_i\times A_{i+1}\to\F$. If
\begin{equation}
                \sum_{i=1}^r g_i(a_i,a_{i+1})=0
                \qquad\text{for all }a\in\prod_i A_i,
 \label{eq:zero-cycle-sum}
\end{equation}
then there are functions $Y_i:A_i\to\F$ with
$g_i(u,v)=Y_i(u)+Y_{i+1}(v)$ for every $i,u,v$.
\end{lemma}

\begin{proof}
Fix a base address $b$. Taking the rectangular difference in
coordinates $i$ and $i+1$ of Equation~\eqref{eq:zero-cycle-sum}
eliminates all terms except $g_i$, because $r\ge3$. Consequently,
with
\begin{align*}
 c_i&=g_i(b_i,b_{i+1}),\\
 \alpha_i(u)&=g_i(u,b_{i+1})+c_i,\\
 \gamma_i(v)&=g_i(b_i,v)+c_i,
\end{align*}
we have $g_i(u,v)=c_i+\alpha_i(u)+\gamma_i(v)$.
Varying coordinate $i$ alone from the base address gives
$\alpha_i(u)=\gamma_{i-1}(u)$. Evaluation at $b$ gives
$\sum_i c_i=0$.

Choose constants $C_i\in\F$ with $C_i+C_{i+1}=c_i$ cyclically:
choose $C_1$ arbitrarily and recurse, the final equation holding
precisely because the $c_i$ sum to zero. Set
$Y_i(u)=\alpha_i(u)+C_i$. Then
\[
 Y_i(u)+Y_{i+1}(v)
   =c_i+\alpha_i(u)+\gamma_i(v)=g_i(u,v).
\]
For a singleton factor the corresponding deviations vanish, so
the same argument includes that case.
\end{proof}

\subsection{The isomorphism criterion at full addresses}

The image sets extracted from the game may be large, and their scalar
equalities do not choose compatible vectors. We now characterize the
stronger condition needed for an isomorphism: one row-sum function for
each real type, satisfying every comparison pointwise. This criterion
applies to arbitrary square shifts, not only the special twists used
later to construct a counterexample to identification.

\begin{proposition}[Global offsets]
\label{prop:global-offsets}
For any right-hand sides $\beta$ in the construction of
Section~\ref{sec:encoding}, all square equations admit a simultaneous
assignment to the readout names if and only if there are row-sum
functions
$h_P:\{0,1\}^r\to\F^{d_P}$, one for each real type, satisfying
\begin{equation}
 \lambda_{e,0}(h_P(a))+\lambda_{e,1}(h_R(\sigma_e(a)))
       =\Delta_e(a)
 \qquad(e:P\to R,\ a\in D_e).
 \label{eq:global-offsets}
\end{equation}
Equivalently, these functions exist exactly when
$G(0)\cong G(\beta)$.
\end{proposition}

\begin{proof}
Given global offsets, sum their horizontal readouts around every
real-type site ring to define $h_P$. These functions have row-sum
form by definition. Summing the square equations around a link
gives Equation~\eqref{eq:link-sum}, now for the assigned offsets
rather than game responses. Summing along the wire cancels the
internal scalar sums and gives
Equation~\eqref{eq:global-offsets}.

Conversely, suppose the displayed functions exist. At each internal
type $s$ of a wire $e:P\to R$, define a scalar function on $D_e$ by
\begin{equation}
 h_s(a)=\lambda_{e,0}(h_P(a))
       +\sum_{\ell\text{ from }P\text{ through }s}\Delta_\ell(a),
 \label{eq:internal-offset}
\end{equation}
where the sum includes the link arriving at $s$. Every $h_s$ has
row-sum form. Along the first and internal links, the link identity
\eqref{eq:link-sum} follows immediately by cancellation. It holds
also on the last link by Equation~\eqref{eq:global-offsets}, with
$h_R$ evaluated at $\sigma_e(a)$. The latter function remains of
row-sum form in $a$ because the transmission acts coordinatewise.

Choose a row decomposition of each real-type function on the full
binary product, and of each internal-type function on its allowed
product. Assign the chosen row terms to all horizontal readout
names of that type. Each choice is made once per type, so readouts
shared between different incident wires already have consistent
values.

Fix a link $\ell=tu$ on wire $e$. For its matched addresses, denote
the two chosen row-$i$ horizontal values by $x_t$ and $x_u$.
The residual
\[
 g_i(a_i,a_{i+1})=
 \beta_{B(i,\ell,a)}+L_{\ell,t}(x_t)+L_{\ell,u}(x_u)
\]
is a function of those two source coordinates, and the link
identity proves $\sum_i g_i(a_i,a_{i+1})=0$ on $D_e$.
Lemma~\ref{lem:cycle-lifting} supplies functions
$Y_i:D_{e,i}\to\F$ whose consecutive sums equal $g_i$.
Assign $Y_i(a_i)$ to the vertical readout of this link indexed by
the match corresponding to $a_i$.

This assignment is well defined: on every link, including the
last, its matches are indexed bijectively by source-domain labels.
Every square uses an adjacent pair of such labels that extends
to a full address because all other product factors are nonempty.
Thus every square equation is satisfied, not just equations from
a selected collection of addresses. Different links have distinct
vertical readout names, so their assignments are independent.
Site blocks impose no additional constraint. We have obtained
the desired global offsets.

The final equivalence with graph isomorphism is
Proposition~\ref{prop:isomorphism}.
\end{proof}

For later use, note that a row-sum function defined on a product
subdomain of $\{0,1\}^r$ extends to a row-sum function on the whole
cube: extend each of its row terms arbitrarily to the corresponding
binary pair domain. This assertion concerns existence of the
extension, not uniqueness of a row decomposition.
